\chapter{哈希表和前缀和}

哈希表解决“快速查历史信息”，前缀和解决“快速计算区间信息”。两者结合后，可以把很多子数组问题从 \complexity{O(n^2)} 降到 \complexity{O(n)}。

两数之和是哈希表入门，但真正要记住的是模型：当前元素需要查询历史中是否存在某个补数。只要查询对象可以作为 key，哈希表就能把线性查找降为平均常数查找。暴力方法枚举 \code{i} 和 \code{j}，判断 \code{nums[i] + nums[j] == target}，时间复杂度 \complexity{O(n^2)}。慢点不在加法，而在“为了每个数反复查找另一个数”。哈希表版本从左到右扫描，用 \code{value_to_index} 保存已经出现过的值及其下标；处理当前值时，只查 \code{target - nums[i]} 是否在历史中出现过。

\begin{lstlisting}[language=C++,caption={两数之和：保存历史值到下标}]
std::vector<int> two_sum_hash(const std::vector<int>& nums, int target)
{
    std::unordered_map<int, int> value_to_index;
    value_to_index.reserve(nums.size());

    for (int index = 0; index < static_cast<int>(nums.size()); ++index) {
        const int need = target - nums[index];
        const auto found = value_to_index.find(need);
        if (found != value_to_index.end()) {
            return {found->second, index};
        }
        value_to_index.emplace(nums[index], index);
    }

    return {};
}
\end{lstlisting}

这段代码必须先查再插入，避免同一个元素和自己配对。如果数组里有重复值，\code{emplace} 会保留第一次出现的下标；对力扣 1 这种保证唯一答案的题足够。如果题目要求所有答案，就不能这样提前返回，而要继续扫描并处理去重。面试时可以这样表达：哈希表保存历史值到下标的映射，查询 key 是当前值的补数，平均时间 \complexity{O(n)}，空间 \complexity{O(n)}。

前缀和的定义是：\code{prefix[i]} 表示前 \code{i} 个元素的和。那么子数组 \code{[left, right]} 的和是 \code{prefix[right + 1] - prefix[left]}。如果题目问和为 \code{k} 的子数组个数，就等价于对每个当前位置查找之前有多少个前缀和等于 \code{current_prefix - k}。

\begin{lstlisting}[language=C++,caption={和为 K 的子数组}]
int subarray_sum(const std::vector<int>& nums, int k)
{
    std::unordered_map<int, int> prefix_count;
    prefix_count.reserve(nums.size() + 1);
    prefix_count[0] = 1;

    int prefix = 0;
    int answer = 0;
    for (int x : nums) {
        prefix += x;
        const int need = prefix - k;
        if (const auto found = prefix_count.find(need); found != prefix_count.end()) {
            answer += found->second;
        }
        ++prefix_count[prefix];
    }
    return answer;
}
\end{lstlisting}

为什么先查再插入？因为当前前缀只能和之前的前缀组成非空子数组。如果先插入，在 \code{k == 0} 时可能把当前位置和自己配对。为什么保存次数而不是下标？因为题目问个数，一个当前前缀可以和多个历史前缀组成答案。

最长连续序列展示 \code{unordered_set} 的另一种用法。把所有数放进集合后，只从没有前驱 \code{x - 1} 的数开始扩展。这样每条连续链只从头部扫描一次，每个数最多被访问常数次，平均时间 \complexity{O(n)}。

哈希题常见错误包括：用 \code{operator[]} 做纯查询导致插入多余 key；忘记处理重复值；依赖 \code{unordered_map} 遍历顺序；没有考虑 \code{long long} 前缀和溢出；把滑动窗口误用到含负数数组。

最长连续序列要求在未排序数组中找最长的连续整数链。排序可以做，复杂度 \complexity{O(n\log n)}；暴力方法从每个数开始向后找 \code{x + 1}、\code{x + 2}，如果每次都线性查找，会退化到 \complexity{O(n^2)}。哈希集合优化了“某个数是否存在”的查询，但如果从每个数都向后扩展，连续链中间的数会被重复扫描。真正的关键是只从链头开始扩展：如果 \code{x - 1} 存在，说明 \code{x} 不是开头，跳过。

\begin{lstlisting}[language=C++,caption={最长连续序列：只从链头扩展}]
int longest_consecutive(const std::vector<int>& nums)
{
    std::unordered_set<int> values;
    values.reserve(nums.size());
    for (const int x : nums) {
        values.insert(x);
    }

    int best = 0;
    for (const int x : values) {
        if (values.contains(x - 1)) {
            continue;
        }

        int current = x;
        int length = 1;
        while (values.contains(current + 1)) {
            ++current;
            ++length;
        }
        best = std::max(best, length);
    }

    return best;
}
\end{lstlisting}

这题的摊还分析要讲清楚：每条连续链只会从最小值开始扫描一次，链中每个元素只在这次扫描里被访问，因此平均时间是 \complexity{O(n)}。不能在遍历 \code{values} 时删除元素，除非你非常清楚迭代器失效规则；本书这里选择不删除，代码更直接。

四数相加 II 展示“拆半”的哈希思想。题目给四个数组，统计 \code{a + b + c + d == 0} 的元组数量。暴力枚举四个下标是 \complexity{O(n^4)}。如果固定前三个再查第四个，仍然是 \complexity{O(n^3)}。更好的做法是把四个数组拆成两组：先统计所有 \code{a + b} 的和出现次数，再枚举 \code{c + d}，查询 \code{-(c + d)} 出现过多少次。

\begin{lstlisting}[language=C++,caption={四数相加 II：两两拆分}]
int four_sum_count(const std::vector<int>& nums1,
                   const std::vector<int>& nums2,
                   const std::vector<int>& nums3,
                   const std::vector<int>& nums4)
{
    std::unordered_map<int, int> pair_sum_count;
    pair_sum_count.reserve(nums1.size() * nums2.size());

    for (const int a : nums1) {
        for (const int b : nums2) {
            ++pair_sum_count[a + b];
        }
    }

    int answer = 0;
    for (const int c : nums3) {
        for (const int d : nums4) {
            const auto found = pair_sum_count.find(-(c + d));
            if (found != pair_sum_count.end()) {
                answer += found->second;
            }
        }
    }

    return answer;
}
\end{lstlisting}

拆半的本质是把 \complexity{n^4} 的组合空间拆成两个 \complexity{n^2} 的集合，再用哈希表连接它们。空间复杂度是 \complexity{O(n^2)}。这类题的面试表达要说明为什么保存次数：同一个两数和可能由多个下标对产生，每一个都能和后半部分组成不同元组。

\begin{principlebox}{哈希题复盘模板}
哈希题不要只写“用 map”。复盘时必须写：key 是什么，value 是什么；查询发生在插入前还是插入后；重复 key 怎么处理；是否需要次数、最早位置、最新位置还是集合存在性；平均复杂度和最坏情况分别是什么。
\end{principlebox}

本章力扣主线：1 两数之和、49 字母异位词分组、128 最长连续序列、560 和为 K 的子数组、523 连续的子数组和、525 连续数组、454 四数相加 II。复盘时必须说明哈希表保存的是值、下标、次数还是最早位置。
